\documentclass[11pt]{article}
\usepackage[T1]{fontenc}
\usepackage{lmodern}
\usepackage{amsmath,amssymb,amsthm,mathtools}
\usepackage{graphicx}
\usepackage[margin=1in]{geometry}
\usepackage[colorlinks=true,linkcolor=blue,citecolor=blue,urlcolor=blue]{hyperref}
\usepackage{microtype}

\newtheorem{theorem}{Theorem}
\newtheorem{lemma}[theorem]{Lemma}
\newtheorem{corollary}[theorem]{Corollary}
\newtheorem{remark}[theorem]{Remark}
\newcommand{\R}{\mathbb R}
\newcommand{\Sphere}{\mathbb S}
\newcommand{\cHC}{\mathbf c}
\newcommand{\dd}{\,d}

\title{Sharp time decay at the bottom of the spherical spectrum\\
\large A full answer to the first sharpness conjecture in arXiv:2504.21479}
\author{Automated mathematics discovery run \texttt{fa\_banach\_001}\\
Lane 5, candidate for human review}
\date{19 July 2026}

\begin{document}
\maketitle

\begin{center}
\fbox{\parbox{0.92\textwidth}{\textbf{Status: candidate full solution, likely valid.}
The two conjectured pointwise powers are attained with explicit nonzero
leading coefficients.  This packet does not address the separate conjecture
about the $L^1$ growth of the distinguished-Laplacian kernel.}}
\end{center}

\section{Source conjecture}

Kuznetsova and Song study oscillatory functions of shifted Laplacians on a
noncompact symmetric space $G/K$ \cite{KS}.  On page 3, after comparing their
gapless $|t|^{-\nu}$ pointwise estimate with the positive-gap
$|t|^{-\nu/2}$ estimate of Anker and Zhang \cite{AZ}, they write:

\begin{quote}
``Note the difference in the exponent between the two cases, $\nu$ against
$\nu/2$. We conjecture that both are optimal, for small $x^+$, and reflect
different behaviour of the two operators.''
\end{quote}

The complete source context is reproduced below.

\begin{center}
\includegraphics[width=0.98\textwidth]{figures/open_problem_crop.png}
\end{center}

We prove locally uniform asymptotic formulas on every compact spatial set.
In particular, both upper-bound powers are optimal throughout the small
$x^+$ regime.

\section{Setting and statement}

Let $G$ be a connected noncompact semisimple Lie group with finite center,
$K$ a maximal compact subgroup, and $\mathfrak a$ a Cartan subspace of real
dimension $l$.  Let $\Sigma_r^+$ be the set of positive reduced roots,
$d=|\Sigma_r^+|$, and
\[
   \nu=2d+l.
\]
Write $\varphi_\lambda$ for the elementary spherical functions and
$\cHC(\lambda)$ for the Harish-Chandra $c$-function, using the normalization
of \cite{KS}.  Spherical inversion gives the gapless kernel
\begin{equation}\label{eq:gapless-kernel}
 k_t^0(x)=C_G\int_{\mathfrak a^*}e^{it|\lambda|}\psi(|\lambda|)
 \varphi_\lambda(x)|\cHC(\lambda)|^{-2}\dd\lambda .
\end{equation}
Here $C_G\ne0$ is the inversion constant.

For each $\alpha\in\Sigma_r^+$ put
\[
 a_\alpha=\frac{m_\alpha}{4}+\frac12,
 \qquad b_\alpha=\frac{m_\alpha}{4}+\frac{m_{2\alpha}}2,
 \qquad \alpha_0=\frac{\alpha}{\langle\alpha,\alpha\rangle}.
\]
If $c_0$ is the constant in the Gindikin--Karpelevich formula, define
\begin{equation}\label{eq:Ac}
 D=|c_0|^{-2}\prod_{\alpha\in\Sigma_r^+}
       \Gamma(a_\alpha)^2\Gamma(b_\alpha)^2,
 \qquad
 A_{\cHC}=D\int_{\Sphere^{l-1}}
       \prod_{\alpha\in\Sigma_r^+}
       \langle\omega,\alpha_0\rangle^2\dd\omega .
\end{equation}
The integrand is nonnegative and is positive away from finitely many great
subspheres, so $A_{\cHC}>0$.

\begin{theorem}[Sharp gapless and positive-gap asymptotics]\label{thm:main}
Let $Q\subset G$ be compact and let $t\to+\infty$.
\begin{enumerate}
\item If $\psi$ is a Schwartz function on $[0,\infty)$ and
$\psi(0)\ne0$, then, uniformly for $x\in Q$,
\begin{equation}\label{eq:gapless-asymp}
 k_t^0(x)=C_GA_{\cHC}\varphi_0(x)\psi(0)\Gamma(\nu)
 e^{i\pi\nu/2}t^{-\nu}+o_Q(t^{-\nu}).
\end{equation}
\item Let $\varepsilon>0$.  Let $m$ be smooth, compactly supported, even
near zero, and satisfy $m(0)\ne0$.  For
\begin{equation}\label{eq:gapped-kernel}
 k_{t,\varepsilon}(x)=C_G\int_{\mathfrak a^*}
 e^{it\sqrt{|\lambda|^2+\varepsilon}}m(|\lambda|)
 \varphi_\lambda(x)|\cHC(\lambda)|^{-2}\dd\lambda,
\end{equation}
one has, uniformly for $x\in Q$,
\begin{align}\label{eq:gapped-asymp}
 k_{t,\varepsilon}(x)
 &=C_GA_{\cHC}\varphi_0(x)m(0)
 2^{\nu/2-1}\varepsilon^{\nu/4}\Gamma(\nu/2)\notag\\
 &\quad\times e^{it\sqrt\varepsilon+i\pi\nu/4}t^{-\nu/2}
 +o_Q(t^{-\nu/2}).
\end{align}
\item Formula \eqref{eq:gapped-asymp} remains valid for the uncut Bessel
amplitude
\[
       m(r)=(r^2+\varepsilon)^{-\alpha/2},\qquad \operatorname{Re}\alpha>n,
\]
where $n=\dim(G/K)$; in this case $m(0)=\varepsilon^{-\alpha/2}$.
\end{enumerate}
\end{theorem}

Since $\varphi_0$ is positive and has a positive minimum on $Q$, both
leading coefficients are bounded away from zero in modulus on $Q$.
Consequently, for all sufficiently large $t$,
\[
 \inf_{x\in Q}|k_t^0(x)|\ge c_Qt^{-\nu},\qquad
 \inf_{x\in Q}|k_{t,\varepsilon}(x)|\ge c_{Q,\varepsilon}t^{-\nu/2}.
\]
Thus neither power can be uniformly improved on any compact spatial region.

\section{Proof intuition}

The time powers come entirely from the bottom of the spherical spectrum.
Every reduced positive root contributes an exact squared linear zero to the
Plancherel density.  There are $d$ such roots; together with the
$l$-dimensional polar Jacobian this produces a radial amplitude of exact
order $r^{2d+l-1}=r^{\nu-1}$.  With no spectral gap, the phase is $tr$, so
the half-line endpoint Fourier lemma gives $t^{-\nu}$.

With a positive gap, the phase has a quadratic stationary endpoint:
$\sqrt{r^2+\varepsilon}=\sqrt\varepsilon+O(r^2)$.  The exact variable
$s=\sqrt{r^2+\varepsilon}-\sqrt\varepsilon$ converts it to a linear phase,
and the transformed amplitude has order $s^{\nu/2-1}$.  This gives
$t^{-\nu/2}$.  The angular coefficient is an integral of a nonzero square,
so neither leading term can cancel.

\section{Proof}

\subsection{Exact order of the Plancherel density}

The Gindikin--Karpelevich product in \cite[Eq.~(2.6)]{KS} is
\[
 \cHC(\lambda)=c_0\prod_{\alpha\in\Sigma_r^+}
 \frac{2^{-iu_\alpha}\Gamma(iu_\alpha)}
 {\Gamma(a_\alpha+iu_\alpha/2)
  \Gamma(b_\alpha+iu_\alpha/2)},
 \qquad u_\alpha=\langle\lambda,\alpha_0\rangle.
\]
For real $u$,
\[
 \frac{|\Gamma(a+iu/2)|^2|\Gamma(b+iu/2)|^2}
      {|\Gamma(iu)|^2}
 =\Gamma(a)^2\Gamma(b)^2u^2h_{a,b}(u),
\]
where $h_{a,b}$ is positive, even, and real analytic near zero, with
$h_{a,b}(0)=1$.  This follows from $\Gamma(iu)^{-1}=iu+O(u^2)$; evenness
removes the odd term after taking modulus square.  Hence, uniformly for
$\omega\in\Sphere^{l-1}$,
\begin{equation}\label{eq:c-leading}
 |\cHC(r\omega)|^{-2}
 =Dr^{2d}\prod_{\alpha\in\Sigma_r^+}
 \langle\omega,\alpha_0\rangle^2\bigl(1+O(r^2)\bigr).
\end{equation}
This identity remains uniform at root hyperplanes because the squared root
forms have been factored out exactly.

For $r\ge0$ define the polar density
\[
 q_x(r)=r^{l-1}\int_{\Sphere^{l-1}}
 \varphi_{r\omega}(x)|\cHC(r\omega)|^{-2}\dd\omega.
\]
Spherical functions are analytic in $\lambda$ and
$\varphi_{r\omega}(x)=\varphi_0(x)+O_Q(r)$ uniformly for $x\in Q$.
The term linear in $r$ is odd in $\omega$, whereas the leading product in
\eqref{eq:c-leading} is even.  Its angular integral vanishes.  Therefore
\begin{equation}\label{eq:q-leading}
 q_x(r)=A_{\cHC}\varphi_0(x)r^{\nu-1}+O_Q(r^{\nu+1}).
\end{equation}
More precisely, $q_x(r)=r^{\nu-1}H_x(r^2)$ near zero, with $H_x$ smooth
and $H_x(0)=A_{\cHC}\varphi_0(x)$, locally uniformly in $x$.

\subsection{Endpoint Fourier lemma}

We use the standard half-line endpoint formula: if $\mu>0$ and
$g(s)=s^{\mu-1}G(s)$ is compactly supported on $[0,\infty)$, with $G$
smooth near zero, then
\begin{equation}\label{eq:watson}
 \int_0^\infty e^{its}g(s)\dd s
 =G(0)e^{i\pi\mu/2}\Gamma(\mu)t^{-\mu}+o(t^{-\mu}).
\end{equation}
The same conclusion holds after adding a tail that is $O(t^{-N})$ for some
$N>\mu$.  One proof inserts an Abel factor, evaluates
$\int_0^\infty e^{-(\eta-it)s}s^{\mu-1}\dd s
=\Gamma(\mu)(\eta-it)^{-\mu}$, subtracts the leading Taylor term of $G$,
and integrates the remainder by parts.  This also proves uniformity when
$G$ ranges over a bounded family of smooth functions.

\subsection{Gapless asymptotic}

Polar coordinates in \eqref{eq:gapless-kernel} give
\[
 k_t^0(x)=C_G\int_0^\infty e^{itr}\psi(r)q_x(r)\dd r.
\]
By \eqref{eq:q-leading}, the endpoint amplitude is
$A_{\cHC}\varphi_0(x)\psi(0)r^{\nu-1}+O_Q(r^\nu)$.
Apply \eqref{eq:watson} with $\mu=\nu$.  Away from zero the phase derivative
is one.  The derivative bounds for $\cHC^{-1}$ and $\varphi_\lambda$ stated
in \cite[Sec.~2.2]{KS} give polynomial growth for all derivatives of $q_x$,
uniformly on $Q$; the Schwartz decay of $\psi$ permits arbitrarily many
integrations by parts.  This proves \eqref{eq:gapless-asymp}.

\subsection{Positive-gap asymptotic}

Set $a=\sqrt\varepsilon$ and
\[
       s=\sqrt{r^2+\varepsilon}-a.
\]
Then
\begin{equation}\label{eq:substitution}
 r^2=s(s+2a),\qquad \frac{dr}{ds}=\frac{s+a}{r}.
\end{equation}
Using \eqref{eq:q-leading}, the transformed leading amplitude is
\begin{align*}
 q_x(r)m(r)\frac{dr}{ds}
 &=A_{\cHC}\varphi_0(x)m(0)
 (s+a)\{s(s+2a)\}^{(\nu-2)/2}+o_Q(s^{\nu/2-1})\\
 &=A_{\cHC}\varphi_0(x)m(0)
 2^{\nu/2-1}\varepsilon^{\nu/4}s^{\nu/2-1}
 +o_Q(s^{\nu/2-1}).
\end{align*}
The phase is now $e^{ita}e^{its}$.  Formula \eqref{eq:watson} with
$\mu=\nu/2$ proves \eqref{eq:gapped-asymp} for compactly supported $m$.

For $m(r)=(r^2+\varepsilon)^{-\alpha/2}$, split with a cutoff supported near
zero.  The local calculation is unchanged and $m(0)=\varepsilon^{-\alpha/2}$.
On the complementary region, $r/\sqrt{r^2+\varepsilon}$ is bounded below.
The derivative estimate for $\cHC^{-1}$ in \cite[Theorem~2.4]{KS}, together
with the spherical-function derivative estimate, gives for every fixed $j$
\[
      \partial_r^jq_x(r)=O_Q((1+r)^{n-1}).
\]
All derivatives of the Bessel amplitude are bounded by a constant times
$(1+r)^{-\operatorname{Re}\alpha}$ at worst.  Thus every term arising from
repeated nonstationary integration by parts is integrable when
$\operatorname{Re}\alpha>n$.  The tail is $O_Q(t^{-N})$ for every $N$, and
the proof is complete. \qed

\section{Verification, scope, and novelty}

The algebraic checker in
\texttt{code/check\_wave\_endpoint\_constants.py} verifies exact model
Fourier integrals, the coefficient in \eqref{eq:substitution}, and positivity
for an $A_2$ root-system model.  These checks do not replace the proof.

This theorem proves pointwise kernel sharpness.  It does not assert that every
finite-$p$ dispersive operator norm has a matching lower bound.  The condition
$\psi(0)\ne0$ is essential to the stated exponent: if the amplitude vanishes
at zero, the first nonzero Taylor term gives a higher endpoint order.  The
separate conjecture in \cite{KS} concerning the optimal $L^1$ growth exponent
$(\nu-1)/2$ for the distinguished Laplacian remains open.

Before promotion, the run's cheap indexes were searched by arXiv id, title,
exact conjecture sentence, and the terms ``wave kernel'', ``spectral gap'',
and ``sharp time decay''.  A bounded primary arXiv search covered
arXiv:2504.21479, the exact phrase, the authors, and the Anker--Zhang
antecedent arXiv:2010.08467.  No later paper claiming this endpoint asymptotic
or resolving the conjecture was found.  This is not an exhaustive novelty
certification.

\paragraph{Human review recommendation.}
Send to a specialist in spherical harmonic analysis.  The two highest-value
checks are the uniform analytic factorization in \eqref{eq:c-leading} across
root hyperplanes and the stated $\operatorname{Re}\alpha>n$ tail threshold.

\begin{thebibliography}{9}
\bibitem{KS}
Y. Kuznetsova and Z. Song,
\emph{Shifted wave equation on noncompact symmetric spaces},
arXiv:2504.21479v2, 2026.

\bibitem{AZ}
J.-Ph. Anker and H.-W. Zhang,
\emph{Wave equation on general noncompact symmetric spaces},
American Journal of Mathematics 146 (2024), 983--1031;
arXiv:2010.08467.
\end{thebibliography}

\end{document}
